%==============================================================================
\appsection{MCMC Implementation Details}
\label{app:mcmc}
%==============================================================================

\subsection{Sampling Algorithm}

The model is estimated using a block-wise Metropolis-Hastings within Gibbs
sampler. The full parameter vector contains
\py{Q.get('n_total_params', 'N/A')} parameters: 4 per subject ($r$, $\ln\phi$,
$\ln\eta$, $\ln\lambda$) for \py{Q.get('n_subjects', 'N/A')} subjects, plus 8
hyperparameters (a population mean and log-standard-deviation for each of the 4
parameters).

These parameters are divided into 115 blocks:
\py{Q.get('n_subjects', 'N/A')} subject blocks, each containing the 4
parameters for a single subject, and 4 hyperparameter blocks, each containing
the population mean and log-standard-deviation for one parameter (e.g.,
$\mu_r$ and $\ln\sigma_r$).

\noindent Each iteration of the sampler cycles through all 115 blocks, proposing
updates to each block conditional on the current values of all other blocks.
Subject blocks are updated via standard Metropolis-Hastings, while
hyperparameter blocks use the coupled transform sampler described below.

\subsection{Non-Centered Parameterization}

The model uses a non-centered parameterization (NCP) to improve MCMC mixing:
\begin{align}
  \varepsilon_i &\sim N(0, 1) \\
  \theta_i &= \mu + \sigma \cdot \varepsilon_i
\end{align}
The sampler stores $\varepsilon_i$ rather than $\theta_i$. Since the
$\varepsilon_i$ are \textit{a priori} independent of $(\mu, \sigma)$, the
funnel geometry is eliminated in the transformed space.

\subsection{Coupled Transform Sampler}

When proposing $(\mu', \sigma')$, the epsilons are updated as:
\begin{equation}
  \varepsilon_i' = \frac{\mu + \sigma \varepsilon_i - \mu'}{\sigma'}
\end{equation}
Since $\theta_i' = \mu' + \sigma' \varepsilon_i' = \theta_i$ by construction,
the data likelihood cancels from the acceptance ratio:
\begin{equation}
  \log \alpha =
    \underbrace{\log q(\mu, \sigma \mid \mu', \sigma')
    - \log q(\mu', \sigma' \mid \mu, \sigma)}_{\text{proposal ratio}}
    + \underbrace{N(\log\sigma - \log\sigma')}_{\text{Jacobian}}
    + \underbrace{-\tfrac{1}{2}\textstyle\sum_{i=1}^N
      (\varepsilon_i'^2 - \varepsilon_i^2)}_{\text{epsilon prior ratio}}
    + \underbrace{\log p(\mu', \sigma')
    - \log p(\mu, \sigma)}_{\text{hyperprior ratio}}
\end{equation}

\subsection{Likelihood Function}

The likelihood for subject $i$'s choice $y_{it} \in \{A, B\}$ uses the Fechner
(logistic) error specification from (\ref{eq:likelihood}):
\begin{equation}
    P(y_{it} = A | \theta_i) = \frac{1}{1 + \exp\left(-\Delta_{it}\right)}
\end{equation}
where $\Delta_{it} = (U_A - U_B) / [\lambda_i (u(\bar{x}) - u(\underbar{x}))]$.

\subsection{Run Configuration}

The sampler runs \py{Q.get('n_chains', 'N/A')} parallel chains through 6
sequential runs: 1 initial reset run followed by 5 resume runs. Each run
collects 100 iterations with thinning factor 10. The first half is discarded
as burn-in, leaving \py{total_samples} posterior samples for inference.

\subsection{Convergence Diagnostics}

Across all \py{Q.get('n_total_params', 'N/A')} parameters, the maximum
$\hat{R}$ is \py{fmt(Q.get('rhat_max', 0), 3)}, indicating excellent
convergence.

\subsection{Software}

Estimation uses the \texttt{bamcmc} package, a JAX-based MCMC sampler with
GPU acceleration, running \py{Q.get('n_chains', 'N/A')} chains in parallel on
an NVIDIA RTX 4070 GPU with 64-bit floating-point precision.


\pagebreak

